\section{Mate selection today: stated, chosen, reproduced}\label{sec:today}

\noindent A present-day population cannot show what happened in the Pleistocene. It is, however, the only place where every
link of Section~\ref{sec:decomp} can be measured, and it shows which features of human mate selection persist when fertility
is controlled and which do not. Three bodies of evidence need to be kept apart: what people say they want, whom they choose,
and who has children.

\subsection{What people say they want}

Across 37 cultures, and again across 45 countries three decades later, women and men both rank kindness, intelligence and
health near the top; women give more weight to financial prospects and men to youth and attractiveness~\cite{buss1989,walter2020}.
When people must build a partner from a limited budget, physical attractiveness behaves as a necessity for men, status and
resources as necessities for women, and kindness and intelligence as necessities for both~\cite{li2002}; the ordering holds in
Eastern and Western samples~\cite{thomas2020}. Factor analyses of what people want in a partner return three dimensions:
warmth and trustworthiness, vitality and attractiveness, and status and resources~\cite{fletcher1999}; compatibility has
recently been proposed as a fourth, valued most in long-term contexts~\cite{jonason2026}.

\subsection{Whom people choose}

Stated preferences are a weak guide to choices. In speed dating, what people said they wanted did not predict whom they
picked~\cite{todd2007}. A meta-analysis of 97 studies of live and ongoing relationships found that physical attractiveness
predicted romantic evaluations at $r\approx0.40$ and earning prospects at $r\approx0.10$ in \emph{both} sexes, with sex
differences that were small and not significant~\cite{eastwick2014}. Machine-learning models given more than a hundred
self-reported traits and preferences predicted who was generally desired and who was generally desiring, but nothing about
which particular pairs clicked~\cite{joel2017}. In a registered report with 10,358 participants in 43 countries, the match
between a person's stated ideals and a partner's traits predicted how that partner was evaluated, but modestly
($\beta=0.19$ for the pattern across 35 traits, 0.11 beside a normative pattern, 0.04 on average for single traits); both
sexes understated how much attractiveness mattered to them, and women overstated, and men understated, how much earning
potential did~\cite{eastwick2025}. Across 11,196 couples, partners' personalities and traits added nothing to the
prediction of relationship quality beyond each person's own view of the relationship~\cite{joel2020}. The evidence is not
all one way: in a preregistered prospective study, preferences recorded while single did predict the characteristics of later
partners~\cite{gerlach2019}.

Several regularities do hold in behavior. Women in speed dating weighted intelligence more and men attractiveness more, and
men did not value a woman's intelligence or ambition once it exceeded their own~\cite{fisman2006}. In a field experiment that
sent invitations from constructed profiles on a large dating site, both sexes preferred attractive profiles; highly
educated women preferred highly educated men, and highly educated men preferred less educated women to the same degree,
so the two groups most likely to stay single were unattractive men with little education and unattractive women with
much~\cite{egebark2021}. Users of dating sites pursue partners more desirable than themselves~\cite{bruch2018}, yet stated
preferences alone, without search frictions, reproduce the sorting seen in marriages~\cite{hitsch2010}. Women's greater
selectivity is real but conditional: it grows with the number of candidates~\cite{fisman2006} and disappears when women are the
ones who approach~\cite{finkel2009}.

Spouses resemble each other. Partner correlations are positive across 22 traits~\cite{horwitz2023}, strongest for
education, attitudes and substance use; assortment on education and height leaves a detectable imprint on the
genome~\cite{yengo2018,robinson2017}. Educational homogamy in the United States rose from 1960 to 2003~\cite{schwartz2005}. In the
survey analyzed below the correlation between spouses' years of schooling was \GssHomOld{} in 1972--1989 and \GssHomNew{} in
2008--2024. As women's schooling has overtaken men's in most countries, the old pattern of women marrying
up in education has ended~\cite{esteve2016}.

\subsection{Said and done: \SdRows{} four-minute meetings}\label{sec:speed}

\emph{Data.} The Columbia speed-dating experiment~\cite{fisman2006} recruited \SdN{} graduate and professional students
(\SdW{} women, \SdM{} men; mean age \SdAge) to \SdSessions{} sessions in 2002--2004. Each pair talked for four minutes.
Each person then rated the other from 1 to 10 on attractiveness, sincerity, intelligence, fun, ambition and shared
interests, said whether they wished to meet again, and guessed how likely the other was to say yes. Before the event each
had divided 100 points among the six attributes. The public file has one row per person per meeting. We verified it by
checksum against four independent mirrors and reproduce the sample sizes of the original report. Variables describing race
and religion were not read.

\emph{Agreement.} Ratings were centered on each rater's mean, which removes leniency and differences among sessions.
Agreement between two observers of the same person, the one-way intraclass correlation, is \JurOneAtt{} for looks,
\JurOneFun{} for fun, \JurOneInt{} for intelligence, \JurOneAmb{} for ambition and \JurOneSin{} for sincerity
(Table~\ref{tab:speed}). Correlating the means of two disjoint random panels of $k$ raters reproduces the Spearman--Brown
curve from $k=1$ to 8 with a largest gap of \JurGapMax{} (Fig.~\ref{fig:main}B). The law is a statistical necessity when
errors are independent, so it should hold in any population. Independence is the condition that matters. Let $r_1$ be one
observer's squared correlation with the truth. If part of the error is common to all observers, so that their errors
correlate $c$, the error of a panel mean has variance $(1-r_1)[c+(1-c)/k]$, and the accuracy of the pooled verdict tends to
\begin{equation}
r_\infty=\frac{r_1}{r_1+c\,(1-r_1)}\label{eq:ceiling}
\end{equation}
however many observers are added: \JurCeil{} for a trait read as poorly as intelligence is here, at $c=0.3$. Shared error
also raises agreement, so consensus then overstates accuracy. In the experiment every rater met the person separately,
which is why the law fits; people who pass on one another's reports add voices and not observers. What a camp supplies is
many independent witnesses of many occasions. The single-observer agreements, unlike the law, need not travel, and here
they are lowered for intelligence by the narrow range of a sample selected for academic ability. Agreement does not differ by the
sex of the person rated for looks (\JurAttMen{} for men, \JurAttWomen{} for women) and is somewhat higher for men's fun
than for women's (\JurFunMen{} against \JurFunWomen).

\emph{Said and did.} Table~\ref{tab:speed} gives the points allocated before the event and the change in the probability
of a yes per standard deviation of each rating, from a linear probability model with a fixed effect for each chooser. The
model was fitted with the chooser's own ratings and again with the mean rating the partner received from everyone else,
which removes the chooser's halo. Either way looks come first for both sexes, then fun; intelligence earns one to four
points; sincerity and ambition earn none. In the sessions that asked, the day after, what had mattered in their decisions
($n=\SdAftN$), women put looks at \SdAftAttW\% and men at \SdAftAttM\%, well above what either had said beforehand, so the
gap between saying and doing is partly visible to the people themselves.

\emph{Legibility and weight.} Across the five attributes of a person, the effect of the panel's rating on a yes
correlates \SdLegPear{} with the agreement between two observers (rank correlation \SdLegSpear). Five points cannot
separate two readings: that people act on what they can read, or that what is easy to read happens to be what matters for
a first meeting. Prediction P27 separates them by lengthening the encounter and adding observers.

\emph{Access and misreading.} Women said yes to \SdYesW\% of partners and men to \SdYesM\%; \SdMatch\% of meetings ended
in a mutual yes. The share of partners who said yes to a person has a coefficient of variation of \SdPopCVM{} among men
and \SdPopCVW{} among women. Rescaling the 1--10 guess to 0--1, each sex expected a yes from \SdExpW\% of partners; women
received \SdActW\% and men \SdActM\%. The scale is arbitrary, but the sexes gave the same answer and met different
outcomes: men's expectation exceeded their outcome by \SdGapM{} points (\SdGapMCI), women's by \SdGapW{} (\SdGapWCI), a
difference of \SdGapDiff{} (\SdGapDiffCI). Men's overestimate of women's interest is what error-management theory predicts~\cite{haselton2000}. A speed-meeting
study designed to test it found the same in men and an underestimate in women~\cite{perilloux2012}, which these data do
not show. Within a person the guess correlates \SdAccM{} (men) and \SdAccW{} (women) with what the
partner decided.

\emph{Checks.} Self-ratings made before the event correlate \SdSelfAtt{} with the panel's verdict for looks, \SdSelfFun{}
for fun, \SdSelfAmb{} for ambition, \SdSelfInt{} for intelligence and \SdSelfSin{} for sincerity. The panel's rating of
intelligence correlates \SdSat{} (\SdSatCI; $n=\SdSatN$) with the median entrance-test score of the person's undergraduate
college, and its rating of looks \SdSatAtt. A speed date is a staged test in the sense of
Proposition~\ref{prop:staging}: everyone knows they are being assessed. The attribute closest to delivery, sincerity, is
the one that observers agree on least, that self-report tracks least, and that earns nothing.

\subsection{Desire, trust, and how relationships begin}\label{sec:desire}

Three further bodies of work bear on what is assessed. \emph{Desire.} Across 211 studies men report a stronger sex drive
than women ($g=0.69$; 0.54 after adjustment for biased responding), a difference between overlapping
distributions~\cite{frankenbach2022}. Willingness to accept a casual offer depends less on the sex of the person asked than on
who is asking: perceived sexual skill of the proposer predicted acceptance in both sexes, and the large sex difference of the
classic studies disappeared for some proposers~\cite{conley2011}. Within relationships, a partner's responsiveness raises
desire, more strongly in women, by making the other feel valued~\cite{birnbaum2016}. \emph{Trust.} What best predicts the
quality of a relationship is each person's perception of the other's commitment and appreciation; the partner's traits add
nothing further~\cite{joel2020}. Such perceptions are partly accurate and partly projection: in 107 couples, judgments of a
partner's willingness to increase intimacy tracked the partner's own report, while judgments of the partner's aversion to
deception did not and mainly reflected assumed similarity~\cite{wickham2020}. \emph{Beginnings.} In seven samples
($n=1{,}897$) about two-thirds of romantic relationships began as friendships, a route that research on initial attraction
between strangers largely overlooks~\cite{stinson2022}. Together these findings say that people assess two things, what a
partner can do and whether the partner will do it for them, and that the second is learned from conduct over time and not
from a first impression (Proposition~\ref{prop:delivery}).

The same two questions organize person perception in general. Perceivers judge others on warmth, read as intent, and on
competence, read as the ability to act on it~\cite{fiske2007}; moral character counts for more than warmth in the
impressions they form~\cite{goodwin2014}. The two are inferred by different rules. Models of impression formation assume a
negativity bias for morality and a positivity bias for competence, and experiments bear this out, the first most strongly
for extreme behavior and the second for moderate behavior~\cite{wojciszke1993,skowronski1987}. The logic is that of a
restrictive schema~\cite{reeder1979}: one skilled act is evidence of ability, since the unskilled cannot produce it, and
one dishonest act is evidence of character, since the honest do not. Occupational
psychology draws the same line between maximum and typical performance, which correlate about 0.4 in a meta-analysis, with
ability more strongly related to the first~\cite{beus2012,sackett1988}. Proposition~\ref{prop:staging} derives these rules
from what a recognized test can and cannot show. It also accounts for three features of courtship that preference lists
leave unexplained: that it takes time, which models treat as information bought by waiting~\cite{wachtmeister1999,seymour2009};
that women tend to underestimate men's commitment~\cite{haselton2000}; and that rivals disparage one another to prospective
partners~\cite{buss1990}, which supplies what a candidate would withhold.

\subsection{Context changes the weights, not the links}

What is weighted varies with circumstance. On tight budgets, Eastern and Western participants allocated differently for
almost every trait (average $d=0.42$) while still putting the same traits first~\cite{thomas2020}. Compatibility is weighted
more in long-term than in short-term contexts~\cite{jonason2026}. The sign of women's access component has reversed in
countries where earning and childbearing have been made compatible (Section~\ref{sec:gss}). A survey of 117,293 people in 175
countries is now public and allows such moderators to be tested directly~\cite{kowal2025}. Two limits apply. The decomposition
predicts that weights follow measurable conditions, such as insecurity, the burden of care and the availability of partners,
and not ancestry: sorting by group on a dating site reflects opportunity and expectation as well as preference, and users
who receive and answer a message across a racial boundary go on to start more such exchanges themselves~\cite{lewis2013}. And
this paper's own analyses compare no racial, ethnic or national groups, and nothing reviewed here supports ranking any.

Table~\ref{tab:culture} places the comparative evidence against the links. Its rows are published comparisons of
institutions and conditions, and each maps to a parameter of the models: who assesses, the share of occasions that are
unstaged, the enforced share of a candidate's product, the share of one sex left unpaired, the number of domains on which
everyone is compared, and the weight a partner's capacity carries. Three cautions apply to them. They are associations
across societies or communities, with the usual confounding. Stated preferences are the outcome in several, and stated
preferences are weak guides to choices (Section~\ref{sec:speed}). And where two large studies disagree, as on whether sex
differences in stated preferences shrink with gender equality~\cite{zentner2012,walter2020}, we report both. What the
table adds is a way to say what a cultural difference would have to change for selection to change: a society that moves
partner choice from kin to strangers, for example, alters the validity of delivery and leaves the validity of capacity
alone, and that is a testable statement (prediction P28).

\subsection{Who reproduces: an access--number decomposition}\label{sec:gss}

\emph{Data and method.} The General Social Survey has interviewed probability samples of United States adults since
1972~\cite{gss2024}. We use all \GssN{} respondents aged 45 to 79 born 1900--1979 with a known number of children, and four traits:
years of schooling, a ten-word vocabulary test, the prestige score of the respondent's occupation, and (for ages 45--64) the
logarithm of personal income. Each trait is standardized, and fitness made relative, within sex and decade of birth, so that
secular trends in schooling and family size do not enter. Eq.~\eqref{eq:split} then divides the selection differential into
access and number. Survey weights are applied; intervals are percentile bootstraps (\GssBoot{} resamples of respondents within
cells). The code reads the public file and reproduces every number here (Section~\ref{app:repro}).

\emph{The number component is negative for every trait in both sexes} (Table~\ref{tab:gss}, Fig.~\ref{fig:today}B). Among people who
become parents, those with more schooling, larger vocabularies, more prestigious occupations and higher incomes have fewer
children. Timing accounts for much of it. Each year's delay in the first birth goes with \GssTimSlopeW\% fewer children for
mothers and \GssTimSlopeM\% for fathers, schooling correlates \GssAfbRW{} (women) and \GssAfbRM{} (men) with age at first birth,
and the path through timing accounts for \GssTimShareW\% of the schooling--number relation in women and \GssTimShareM\% in men.

\emph{The access component has opposite signs in the two sexes.} For income it is $\GssIncMAcc$ (\GssIncMAccCI) in men and
$\GssIncWAcc$ (\GssIncWAccCI) in women; the difference, \GssIncDiffAcc{} (\GssIncDiffAccCI), is the largest contrast in the
table. Of men in the lowest third of income for their cohort, \GssChlMIncLo\% were childless, against \GssChlMIncHi\% in the
highest third; for women the figures run the other way, \GssChlWIncLo\% and \GssChlWIncHi\%. Occupational
prestige shows the same asymmetry at smaller size. For men's income the two components nearly cancel, leaving a total
differential indistinguishable from zero ($\GssIncMTot$; \GssIncMTotCI). A single correlation between status and fertility
would report ``no selection'' where there are two opposed episodes.

\emph{Three episodes, and the shape of the floor.} Access can be divided once more, into pairing (the mean trait of those who
ever married) and conversion (the further difference between parents and the ever-married); the three terms still sum to the
total (Table~\ref{tab:gssthree}). For men every trait is positive at pairing: income \GssIncMPair{} (\GssIncMPairCI), occupational
prestige \GssPreMPair, schooling \GssSchMPair{} and vocabulary \GssVocMPair. Conversion adds nothing for men's income
(\GssIncMConv) and is negative for schooling and vocabulary. For women no trait has a detectable effect at pairing, and
conversion is negative for all four. By fifths of income, men's childlessness is \GssQfMOne, \GssQfMTwo, \GssQfMThree{} percent in
the three lowest and \GssQfMFive{} in the highest, and the share never married falls from \GssNmMOne\% to \GssNmMFive\%
(Fig.~\ref{fig:today}A). The fall is concentrated at the bottom: the lowest-to-middle difference exceeds the middle-to-highest by
\GssFloorM{} percentage points (\GssFloorMCI). Marriage is an imperfect marker of pairing in recent cohorts, in which many
couples cohabit, so the pairing term is conservative for them.

\emph{The pattern is stable.} It appears in every birth cohort from 1900--1919 to 1960--1979, without weights, in the stricter
age window 50--74, and in surveys before and after 1990 (Tables~\ref{tab:gsscohort} and~\ref{tab:gssrobust}).

\emph{It is not new, and it is not peculiar to the United States.} British men with low incomes are more often childless
while women's income is negatively related to family size~\cite{nettle2008}; in European and United States samples the sex
difference in the status--fertility relation is driven by childlessness~\cite{fieder2007,barthold2012,hopcroft2006,weeden2006}; men
with higher incomes in the United States are more likely to marry, less likely to divorce and less likely to be
childless, and the reverse holds for women~\cite{hopcroft2021}. In Swedish registers the gradient of fertility on lifetime income is
strongly positive for men, and for women it has moved from negative to positive across cohorts born 1940--1970~\cite{kolk2023};
in the Nordic countries generally the negative gradient of women's education has vanished while the least educated men
remain the most often childless~\cite{jalovaara2019}. Swedish men with the lowest intelligence-test scores are the least likely
to have a first child~\cite{kolk2019}. What the present analysis adds is the exact split, for four traits and both sexes in one
sample, and its stability over eighty years of birth cohorts. The contrast across countries matters for interpretation: the
sign of the access term for women tracks institutions, since it has reversed where earning and childbearing have been made compatible. The
positive sign for men is found everywhere it has been looked for.

\emph{How much of the variance is the floor?} The opportunity for selection is $I=\GssIM$ in men and \GssIW{} in women, a small
sex difference of the size reported from other United States data~\cite{jokela2010}. Childlessness accounts for \GssShareM\% of it in
men and \GssShareW\% in women (Eq.~\ref{eq:oppsplit}), and the share has risen from \GssShareMOld\% and \GssShareWOld\% in cohorts
born 1920--1939 to \GssShareMNew\% and \GssShareWNew\% in those born 1960--1979 (Fig.~\ref{fig:today}E). As family sizes converge
on two, whether one becomes a parent at all carries a growing part of what variance remains.

\emph{The pairing gate.} Selection at pairing can act only on reproduction that passes through pairing, and the share that
does has fallen (Table~\ref{tab:gssgate}; Fig.~\ref{fig:verify}C). Among men born 1920--1939, \GssGateMOld\% (\GssGateMOldCI)
of those who never married were fathers, and never-married men accounted for \GssOutMOld\% of the children reported by men.
Among men born 1960--1979 the figures are \GssGateMNew\% (\GssGateMNewCI) and \GssOutMNew\%. For women they rose from
\GssGateWOld\% (\GssGateWOldCI) and \GssOutWOld\% to \GssGateWNew\% (\GssGateWNewCI) and \GssOutWNew\%. At the same time
the gate became more selective. The share of men never marrying rose from \GssNevMOld\% to \GssNevMNew\%, and the pairing
component of the differential on men's income rose from \GssPairMOld{} (\GssPairMOldCI) to \GssPairMNew{} (\GssPairMNewCI),
a change of \GssPairMChange{} (\GssPairMChangeCI). For women the pairing component moved from \GssPairWOld{}
(\GssPairWOldCI) to \GssPairWNew{} (\GssPairWNewCI), a change of \GssPairWChange{} (\GssPairWChangeCI): in the most recent
cohort women with higher incomes are slightly more likely to have married, where the pooled estimate of
Table~\ref{tab:gssthree} shows no effect. Event-history data likewise find that women's earnings became more important for
marriage formation between the early and late baby-boom cohorts~\cite{sweeney2002}. Two cautions apply. The survey records
marriage, not partnership, and most births outside marriage now occur to cohabiting couples, with a strong gradient by
education~\cite{lundberg2016}; and men may report children born outside marriage less completely than women do. The
cohort comparison was not planned in advance and should be read as description.

\emph{Standing and sexual partners.} Among respondents aged 45--79 the rank correlation between income and the lifetime
number of sexual partners of the other sex is \GssPartIncM{} (\GssPartIncMCI) in men and \GssPartIncW{} (\GssPartIncWCI)
in women; for vocabulary it is \GssPartVocM{} and \GssPartVocW, for schooling \GssPartSchM{} and \GssPartSchW, and for
occupational standing \GssPartPreM{} and \GssPartPreW{} (medians of \GssPartMedM{} partners for men and \GssPartMedW{} for
women). In this population a man's standing predicts whether he has a partner and whether he marries, not how many
partners he has; the small positive association in women is what later marriage would produce. Reported partner counts
are known to be biased, differently in the two sexes~\cite{mitchell2019}.

\emph{Marital history.} Table~\ref{tab:gsshist} and Fig.~\ref{fig:today}F give family size by marital history, relative
to the mean of the respondent's sex and decade of birth. Against people in a first marriage, remarried men have \GssRemM{} (\GssRemMCI) and
remarried women \GssRemW{} (\GssRemWCI), a sex difference of \GssRemDiff{} (\GssRemDiffCI). Men who divorced and did not
remarry have \GssDivM{} (\GssDivMCI), and \GssDivChlM\% of them are childless against \GssIntChlM\% of men in a first
marriage; for women the figures are \GssDivW{} (\GssDivWCI), \GssDivChlW\% and \GssIntChlW\%. The question about divorce
was put only to the currently married and widowed, and history is not randomly assigned: childless marriages end more
often. The asymmetry agrees with earlier findings that serial marriage raises men's reproductive success and not
women's~\cite{jokela2010}.

\emph{Who marries up.} Among currently married respondents the wife had more schooling than the husband in
\GssHypMoreOld\% of marriages of people born 1920--1939 and less in \GssHypLessOld\%; for those born 1960--1979 the figures
are \GssHypMoreNew\% and \GssHypLessNew\%, and for those born from 1980, \GssHypMoreYoung\% and \GssHypLessYoung\%
(a difference of \GssHypDiffYoung{} points; \GssHypDiffYoungCI; Table~\ref{tab:gsshyp}). The correlation between spouses'
schooling stayed near \GssHypR. In the oldest cohort, born 1900--1919, wives were also the more schooled spouse more often
than not, so husbands' advantage in schooling belongs to the middle of the century and is not a constant that recently
ended. The reversal follows the reversal of the sex gap in schooling and is seen
worldwide~\cite{esteve2016}.

\emph{Genetic evidence points the same way.} Polygenic scores for educational attainment are negatively associated with
number of children in Iceland, the United States and the United Kingdom, largely through later
childbearing~\cite{kong2017,beauchamp2016,hughjones2022,barban2016}; the association reverses once age at first birth is
controlled~\cite{hughjones2022}. On the access side, people who carry rare variants that damage genes under strong selective
constraint are more often childless, more so men than women, apparently because they are less likely to find a partner; the
authors estimate that this accounts for about a fifth of the purifying selection against such variants, while explaining less
than one percent of childlessness~\cite{gardner2022}.

\emph{Limits.} These are phenotypic differentials. They predict genetic change only to the extent that the covariance
between trait and fitness is genetic, which observational data cannot show~\cite{morrissey2010}. Income at ages 45--64 is partly
a consequence of family formation, with opposite effects on mothers' and fathers' earnings~\cite{kleven2019}; schooling is
mostly complete before childbearing and vocabulary is comparatively stable, which is why all four are shown. Respondents
are survivors to at least 45, number of children is top-coded at eight, and everything is self-reported. The estimates
describe one country. They carry no implication for the worth of any person or group and are no basis for policy.

\subsection{Sexual behavior: partners are not offspring}

Sexual selection theory expects reproductive success to rise with the number of mates, more steeply in
males~\cite{jones2009}. In this population it does not rise at all. Among adults aged 45--79, mean family size is highest
for those reporting exactly one lifetime partner of the other sex (\GssBatMOne{} children for men, \GssBatWOne{} for women) and
falls with each further category, to \GssBatMTwenty{} and \GssBatWTwenty{} for twenty or more (Fig.~\ref{fig:today}C). The slope of
relative fitness on relative partner number is \GssBatMGrad{} (\GssBatMGradCI) in men and \GssBatWGrad{} (\GssBatWGradCI) in women.
Reported counts are imperfect: men report more partners than women (medians \GssBatMMed{} and \GssBatWMed{} here), a gap that
estimation and social norms largely explain~\cite{mitchell2019}. The sign of the relation does not depend on such errors. With
effective contraception, status still predicts men's mating opportunities but not their fertility~\cite{perusse1993}, and what
predicts children is the time spent in a stable union: in a natural-fertility population, years married, not number of
partners, carried the relation between mating and reproductive success~\cite{borgerhoffmulder2019}.

Access, by contrast, has been tightening. At ages 22--34 the share reporting no sexual partner in the past year rose from
\GssPlMOld\% to \GssPlMNew\% among men and from \GssPlWOld\% to \GssPlWNew\% among women between 1989--2002 and 2016--2024
(Fig.~\ref{fig:today}D). Among men it is graded by income: \GssPlMLo\% in the lowest third against \GssPlMHi\% in the highest in the
latest period, and \GssPlMNone\% among men with no earnings. An analysis of the same survey for 2000--2018 found the rise
concentrated among unmarried men and associated with low income and with part-time or no employment~\cite{ueda2020}.

\subsection{Eight claims about present-day mate choice, checked}

Table~\ref{tab:audittoday} continues the audit of Table~\ref{tab:audit} for claims about the present.

\begin{table*}[t]
\caption{Eight common claims about present-day mate choice and selection, checked.}\label{tab:audittoday}
\scriptsize
\begin{tabularx}{\textwidth}{@{}P{0.028\textwidth}P{0.20\textwidth}P{0.10\textwidth}L@{}}
\toprule
 & Claim & Verdict & Evidence\\
\midrule
N11 & What people say they want is what they choose. & Weak guide &
Stated preferences did not predict speed-dating choices~\cite{todd2007}; attractiveness and earning prospects predict
evaluations about equally in both sexes whatever people state~\cite{eastwick2014}; pair-specific attraction is not predictable
from traits and preferences~\cite{joel2017}. Against this, preferences recorded while single predicted later partners'
characteristics in one prospective study~\cite{gerlach2019}.\\
\addlinespace
N12 & Women choose on status; men choose on looks. & Half right &
The difference in stated priorities holds across cultures~\cite{buss1989,walter2020,li2002}. In live evaluation both sexes
respond mainly to attractiveness~\cite{eastwick2014,egebark2021}. Where the sexes differ is in outcomes: income raises men's chances
of marriage and fatherhood and lowers women's chances of motherhood (Table~\ref{tab:gss})~\cite{nettle2008,hopcroft2021}, and
education preferences in a field experiment had opposite signs~\cite{egebark2021}.\\
\addlinespace
N13 & Intelligence is attractive. & Perceived, not measured &
Measured intelligence can be detected from videos and in live interaction but does not raise attractiveness; perceived
intelligence and funniness do~\cite{driebe2021}. Humor ability tracks intelligence and predicts mating
success~\cite{greengross2011}. Men did not value women's intelligence above their own~\cite{fisman2006}. General capacity reaches
choice through performance and standing more than through direct attraction.\\
\addlinespace
N14 & Women are the choosy sex. & Conditional &
Women are more selective in most settings, more so with more candidates~\cite{fisman2006}; the difference vanished when women
approached and men sat~\cite{finkel2009}. Human mating does not follow one universal pattern of sex roles~\cite{brown2009}, both
sexes choose~\cite{stewartwilliams2013}, and in most small-scale societies kin choose too~\cite{walker2011}.\\
\addlinespace
N15 & Women's preferences shift toward masculine men when fertile. & Not supported &
A meta-analysis of 58 reports was largely nonsupportive, with effects confined to published studies using imprecise
definitions of the fertile phase~\cite{wood2014}; a large longitudinal study with hormone assays found no relation between
hormonal status and preference for facial masculinity~\cite{jones2018}.\\
\addlinespace
N16 & Masculine faces and deep voices mark the men who reproduce most. & Mostly not &
Across 96 studies, strength and muscularity were the only consistent predictors of both mating and reproduction; voice
pitch, height and testosterone predicted mating only, and facial masculinity predicted neither~\cite{lidborg2022}.\\
\addlinespace
N17 & For men, more partners mean more offspring. & Not now &
Family size peaks at one lifetime partner in both sexes and declines thereafter (Fig.~\ref{fig:today}C). Status predicts men's
mating opportunities but not fertility under contraception~\cite{perusse1993}. Men with three or more successive spouses had 19\%
more children than men with one, with no such effect in women~\cite{jokela2010}: unions, not partners, are what count.\\
\addlinespace
N18 & Selection on humans has stopped; or, intelligence is being bred out. & Neither &
Selection continues on many traits~\cite{byars2010,sanjak2018,gardner2022,akbari2026}. On education-linked variants it is negative
and slow, about 0.01 standard units per decade in Iceland~\cite{kong2017}, too slow to account for more than a small part of
observed change~\cite{beauchamp2016}. Test-score trends in both directions arise within families and are
environmental~\cite{bratsberg2018}; brain volumes rose across twentieth-century birth cohorts~\cite{decarli2024}.\\
\bottomrule
\end{tabularx}
\end{table*}
